\documentclass[a4paper,11pt]{article}
\usepackage[utf8]{inputenc}
\usepackage[T1]{fontenc}
\usepackage[french]{babel}
\usepackage[left=2cm,right=2cm,top=2cm,bottom=2cm]{geometry}
\usepackage{xcolor}
\usepackage{hyperref}
\usepackage{fontawesome5}
\usepackage{parskip}

% --- Couleurs Harmattan AI ---
\definecolor{primary}{RGB}{30, 60, 90}

\hypersetup{
    colorlinks=true,
    urlcolor=primary,
}

\begin{document}

% --- EXPÉDITEUR (haut gauche) ---
\noindent
\begin{minipage}[t]{0.5\textwidth}
    \textbf{Loïc Aron MBASSI EWOLO}\\
    Élève-Ingénieur Bac+5 -- Double diplôme ENSEA\\[3pt]
    \faMapMarker\ Cergy, France\\
    \faPhone\ (+33) 7 59 52 33 98\\
    \faEnvelope\ \href{mailto:loicaron.mbassiewolo@ensea.fr}{loicaron.mbassiewolo@ensea.fr}
\end{minipage}
\hfill
% --- LIEU ET DATE (haut droite) ---
\begin{minipage}[t]{0.4\textwidth}
    \raggedleft
    Cergy, le 27 janvier 2026
\end{minipage}

\vspace{1.2cm}

% --- DESTINATAIRE (à droite) ---
\hfill
\begin{minipage}[t]{0.5\textwidth}
    \raggedleft
    \textbf{Harmattan AI}\\
    Service Recrutement\\
    Paris, France
\end{minipage}

\vspace{1cm}

\noindent\textbf{Objet :} Candidature au poste de FPGA Intern

\vspace{0.8cm}

Madame, Monsieur,

Élève-ingénieur en 5\textsuperscript{e} année à l'\textbf{École Polytechnique de Yaoundé} et actuellement en double diplôme à l'\textbf{ENSEA} en Traitement du Signal et Systèmes Embarqués, je souhaite rejoindre Harmattan AI en tant que FPGA Intern. Votre mission de développement de systèmes de défense autonomes et scalables correspond à mes aspirations professionnelles et à mon parcours technique.

Je suis actuellement impliqué dans le challenge \textbf{CoHoMa}, un projet de robotique autonome militaire organisé par l'armée française en collaboration avec l'ENSEA. Dans ce cadre, je travaille sur des \textbf{Nvidia Jetson} équipées de \textbf{Lidar 3D VLP16} et de \textbf{caméras de profondeur}. J'y développe des algorithmes de \textbf{SLAM} et de fusion de données capteurs pour la navigation autonome. Cette expérience m'a confronté aux exigences des systèmes critiques : fiabilité, temps réel et optimisation des ressources embarquées.

Parallèlement, je suis en train d'approfondir mes compétences en \textbf{FPGA} à l'ENSEA, où je conçois des circuits numériques en \textbf{VHDL}. Bien que cette formation soit en cours, ma solide base en \textbf{traitement du signal} (DSP, filtrage numérique, MATLAB/Simulink) et en programmation bas niveau (\textbf{C/C++}, Linux) me permet d'aborder rapidement les concepts de simulation radar et d'implémentation hardware.

Mon projet académique sur la \textbf{détection de défauts et le contrôle tolérant aux pannes} d'un drone m'a également formé à la modélisation de systèmes dynamiques et à la simulation de scénarios de défaillance -- des compétences directement transférables au benchmarking de systèmes radar.

Je suis convaincu que mon profil hybride -- alliant expérience terrain en systèmes embarqués défense, formation FPGA en cours, et rigueur acquise en recherche (MILA) -- me permettrait de contribuer efficacement au développement de votre simulateur de cibles radar.

Je reste à votre disposition pour un entretien afin de discuter de ma candidature.

Je vous prie d'agréer, Madame, Monsieur, l'expression de mes salutations distinguées.

\vspace{0.8cm}

\textbf{Loïc Aron MBASSI EWOLO}\\
\textit{Élève-Ingénieur ENSEA / Polytechnique Yaoundé}

\end{document}
